% ECONOMICS_PROBLEM: {"id": "G12-270", "title": "Risk versus mispricing in return factors", "jel_code": "G12", "source_id": "OP-0314"}
% FIRST_POSTED: 2026-10-09
% ATTEMPT_STATUS: unresolved
% ENTRY_KIND: research_attempt
\documentclass[11pt,a4paper]{article}
\usepackage[T1]{fontenc}
\usepackage[utf8]{inputenc}
\usepackage{lmodern}
\usepackage[margin=26mm]{geometry}
\usepackage{amsmath,amssymb,amsthm,mathtools}
\usepackage[hidelinks]{hyperref}
\usepackage{microtype}
\setlength{\parskip}{4pt}
\newtheorem{theorem}{Theorem}[section]
\newtheorem{proposition}[theorem]{Proposition}
\newtheorem{lemma}[theorem]{Lemma}
\newtheorem{corollary}[theorem]{Corollary}
\theoremstyle{definition}
\newtheorem{definition}[theorem]{Definition}
\theoremstyle{remark}
\newtheorem{remark}[theorem]{Remark}
\newcommand{\R}{\mathbb R}
\newcommand{\E}{\mathbb E}
\newcommand{\Prb}{\mathbb P}
\newcommand{\ind}{\mathbf 1}
\newcommand{\Var}{\operatorname{Var}}
\newcommand{\supp}{\operatorname{supp}}
\newcommand{\TV}{\operatorname{TV}}
\DeclareMathOperator*{\argmin}{arg\,min}
\DeclareMathOperator*{\argmax}{arg\,max}
\title{Factor pricing as a joint identified set: risk restrictions, costs and validation}
\author{GPT 6 Astra Ultra}
\date{9 October 2026}
\begin{document}
\maketitle
\noindent\textbf{Catalogue problem:} \texttt{G12-270}. \textbf{Status:} Unresolved. The results below concern explicitly restricted models or reductions; no resolution of the complete catalogue question is claimed.

\section{Question and interpretation}
The catalogue asks whether return factors reflect compensation for risk or mispricing, after executable trading costs, under a risk model restricted independently of the tested returns. Replication across samples, as studied by Jensen, Kelly and Pedersen~\cite{jkp}, is relevant to the stability of factor evidence, but does not by itself identify a risk-versus-belief mechanism. Here we solve a finite-state identified-set problem and give a finite-class validation test. No trading recommendation or empirical classification is made.

\section{An externally specified stochastic-discount-factor class}
There are $S$ states with known probabilities $\pi_s>0$, $\sum_s\pi_s=1$. Payoffs and stochastic discount factors are vectors in $\R^S$, with inner product $\langle a,b\rangle=\sum_s\pi_sa_sb_s$ and norm $\|a\|_2$. Normalize the risk-free gross return to one. A predeclared collection $f_1,\ldots,f_d$ of excess payoffs is required to be priced. Let
\[
 \mathcal A=\{M:\E M=1,\ \E[Mf_\ell]=0\ (1\leq\ell\leq d)\}.
\]
Assume this affine set is nonempty. The economically admissible class is
\begin{equation}\label{eq:class}
 \mathcal M=\{M\in\mathcal A:\|M-1\|_2\leq\sigma,
                  \ \underline m\leq M_s\leq\overline m\ (1\leq s\leq S)\},
\end{equation}
where $0<\underline m\leq1\leq\overline m<\infty$ and $\sigma\geq0$ are also fixed independently of the tested returns. It may be empty; this is a specification failure, not evidence of an infinite premium.

For tested gross excess return $r_j$ and nonnegative executable cost $c_j$, define
\[
 \delta_j(M)=\E[Mr_j],\qquad
 \alpha_j(M)=\E[M(r_j-c_j)].
\]
The target is the joint image of $\mathcal M$ under all these linear functionals. The collection of discount factors is not estimated by choosing whichever one makes the tested factor look priced.

\section{A sharp ellipsoidal image when positivity is inactive}
Let $H=\{h:\E h=0,\ \E[hf_\ell]=0\ \forall\ell\}$, and let $M_0$ be the point of $\mathcal A$ closest to the constant vector $1$. Orthogonal projection exists uniquely in finite dimension. For arbitrary predeclared payoff vectors $g_1,\ldots,g_J$, put
\[
 \widetilde g_j=\operatorname{Proj}_H g_j,\quad
 b_j=\E[M_0g_j],\quad
 G_{jk}=\langle\widetilde g_j,\widetilde g_k\rangle,
 \quad R^2=\sigma^2-\|M_0-1\|_2^2.
\]
Use $g_j=r_j$ and $g_{K+j}=r_j-c_j$ for the joint pricing-and-cost target when there are $K$ tested factors.

\begin{theorem}[Exact joint image]\label{th:ellipse}
If $R^2<0$, the class~\eqref{eq:class} is empty. Suppose $R^2\geq0$. Without the coordinate bounds on $M$, the exact image of the affine norm ball is
\begin{equation}\label{eq:ellipse}
 \{b+v:v\in\operatorname{range}(G),\ v^TG^+v\leq R^2\},
\end{equation}
where $G^+$ is the Moore--Penrose inverse, including the convention that $G=0$ gives image $\{b\}$. The same image is exact for $\mathcal M$ if
\begin{equation}\label{eq:inactive}
 M_{0,s}-R/\sqrt{\pi_s}\geq\underline m,
 \qquad M_{0,s}+R/\sqrt{\pi_s}\leq\overline m
 \quad\hbox{for every }s.
\end{equation}
In general~\eqref{eq:ellipse} is only an outer set; the exact image is obtained by retaining those $v$ for which some $h\in H$ satisfies
\[
 \|h\|_2\leq R,\quad
 \langle h,\widetilde g_j\rangle=v_j\ \forall j,\quad
 \underline m\leq M_{0,s}+h_s\leq\overline m\ \forall s.
\]
This is a finite-dimensional convex feasibility problem.
\end{theorem}
\begin{proof}
Every element of $\mathcal A$ is $M_0+h$ for $h\in H$. The first-order condition for the projection of $1$ says $M_0-1$ is orthogonal to $H$, so
\[
 \|M_0+h-1\|_2^2=\|M_0-1\|_2^2+\|h\|_2^2.
\]
This proves emptiness when $R^2<0$ and reduces the norm constraint to $\|h\|_2\leq R$. Its image coordinate is $b_j+\langle h,\widetilde g_j\rangle$.

The linear map $h\mapsto(\langle h,\widetilde g_j\rangle)_j$ has image $\operatorname{range}(G)$: its kernel is the orthogonal complement in $H$ of the span of the $\widetilde g_j$, and its Gram matrix is $G$. More explicitly, for $v\in\operatorname{range}(G)$ the vector
\[
 h_v=\sum_j(G^+v)_j\widetilde g_j
\]
maps to $GG^+v=v$ and has squared norm $v^TG^+v$. Any other vector with the same image differs from $h_v$ by a vector orthogonal to every $\widetilde g_j$, hence orthogonal to $h_v$. It has weakly greater squared norm. Thus an image $v$ is attainable in the radius-$R$ ball exactly when the inequality in~\eqref{eq:ellipse} holds.

For any $\|h\|_2\leq R$, the individual term $\pi_sh_s^2$ is at most $R^2$, so $|h_s|\leq R/\sqrt{\pi_s}$. Condition~\eqref{eq:inactive} makes every such vector coordinate-feasible. Without it, retaining precisely the vectors admitting the displayed feasible preimage imposes the original coordinate restrictions, proving exactness of the final formulation.
\end{proof}

\begin{corollary}[Sharp bounds for one contrast]
For a single payoff $g$, the unconstrained affine norm-ball image is
\[
 \left[\E[M_0g]-R\|\operatorname{Proj}_Hg\|_2,
       \E[M_0g]+R\|\operatorname{Proj}_Hg\|_2\right].
\]
If its two endpoint discount factors
$M_0\pm R\operatorname{Proj}_Hg/\|\operatorname{Proj}_Hg\|_2$
are coordinate-feasible, these endpoints and every intervening value are attainable in $\mathcal M$. If the projected payoff is zero, the functional is constant on every nonempty admissible class.
\end{corollary}
\begin{proof}
Cauchy--Schwarz bounds $\langle h,\operatorname{Proj}_Hg\rangle$ by $R\|\operatorname{Proj}_Hg\|_2$ and the displayed directions attain equality. Convex mixtures of feasible endpoints remain feasible and trace the full interval. A zero projection is orthogonal to every admissible difference $h$, yielding constancy.
\end{proof}

Coordinate intervals for many factors generally do not form the joint identified set: their extremes may require incompatible discount factors. The Gram matrix in Theorem~\ref{th:ellipse} records that dependence. Positivity restrictions must not be discarded merely because the unconstrained formula is convenient.

\section{Costs and a limit to interpreting positive means}
\begin{proposition}[Cost wedge and a two-state rationalization]\label{pr:cost}
For every $M\in\mathcal M$,
\[
 \delta_j(M)-\alpha_j(M)=\E[Mc_j],\qquad
 \underline m\E c_j\leq\delta_j(M)-\alpha_j(M)
                         \leq\overline m\E c_j.
\]
Thus a gross-priced factor with $\delta_j(M)=0$ has net moment $\alpha_j(M)=-\E[Mc_j]\leq0$.

Separately, in a two-state economy with probabilities $p,1-p\in(0,1)$ and excess return $r=(a,-b)$, $a,b>0$, the strictly positive discount factor
\[
 M_+=\frac{b}{p(a+b)},\qquad
 M_-=\frac{a}{(1-p)(a+b)}
\]
satisfies $\E M=1$ and $\E[Mr]=0$, irrespective of the sign of $\E r$.
\end{proposition}
\begin{proof}
Subtract the definitions of $\delta_j$ and $\alpha_j$. Since $c_j\geq0$, multiply the coordinate bounds on $M$ by $\pi_sc_{j,s}$ and sum to obtain the two inequalities. The gross-pricing implication follows immediately. In the two-state example,
$pM_++(1-p)M_-=(a+b)/(a+b)=1$, while
$pM_+a-(1-p)M_-b=ab/(a+b)-ab/(a+b)=0$.
\end{proof}

The two-state discount factor need not satisfy a particular externally specified bound or price other assets. Its purpose is narrower: a positive ordinary mean does not, without risk restrictions, identify mispricing. Conversely rejection of an externally restricted SDF class may reject that risk class without establishing a particular belief distortion. Holdings, consumption or belief measurements are additional restrictions, not consequences of the return moments alone.

\section{A preregistered finite-class validation test}
Suppose a separate development sample has fixed $K$ strategies, their executable costs, and $L_0$ candidate discount factors $M^{(1)},\ldots,M^{(L_0)}$. On an untouched sample of $T$ independent identically distributed states, assume $0<M^{(\ell)}\leq U$ and $|r_j|\leq R_0$. Let
\[
 \widehat\delta_{\ell j}=T^{-1}\sum_{t=1}^T M_t^{(\ell)}r_{j,t},\qquad
 b_T=UR_0\sqrt{\frac{2\log(2L_0K/\alpha)}{T}}.
\]
The candidate functions are fixed independently of these validation observations.

\begin{theorem}[Simultaneous finite-class pricing test]
With probability at least $1-\alpha$, all $L_0K$ population pricing moments lie in $[\widehat\delta_{\ell j}-b_T,\widehat\delta_{\ell j}+b_T]$. Reject the assertion that some candidate prices all tested gross factors only when, for every $\ell$, at least one of its $K$ intervals excludes zero. If some candidate actually prices all factors, this rejection probability is at most $\alpha$. Replacing $r_j$ by $r_j-c_j$, with $0\leq c_j\leq C_0$, gives simultaneous net-moment intervals using $U(R_0+C_0)$ in place of $UR_0$.
\end{theorem}
\begin{proof}
Each summand is in $[-UR_0,UR_0]$. Hoeffding's inequality gives failure probability at most $\alpha/(L_0K)$ for each displayed interval. A union bound proves simultaneous coverage. On that event, any candidate whose $K$ population moments are all zero has zero in all its intervals, so the stated rule cannot reject. The net-return summands satisfy the asserted absolute bound, giving the last claim by the same argument.
\end{proof}

The test handles the declared finite search and sample splitting only. Serial dependence, an infinite estimated SDF class, costs estimated from overlapping data, or strategies reselected after validation require additional analysis. These limits, together with the absence of a belief or holdings model, leave the original economic decomposition unresolved. The established output is a sharp restricted identified set, exact cost accounting and an honest finite-class test.

\begin{thebibliography}{9}
\bibitem{jkp} T. I. Jensen, B. Kelly and L. H. Pedersen (2023), ``Is There a Replication Crisis in Finance?,'' \emph{The Journal of Finance}. \url{https://doi.org/10.1111/jofi.13249}. Primary publisher record and author publication list: \url{https://www.bryankellyacademic.org/}. This reference concerns factor replication; no risk-versus-mispricing identification theorem is attributed to it.
\end{thebibliography}

\end{document}
